\chapter{Méthodologie de détection des signaux}

\section{Positionnement général de la méthode}
La couche signal met en oeuvre une chaîne de recherche appliquée dont l'objectif
est de transformer une famille d'indicateurs techniques en un signal final
robuste, interprétable et réutilisable. Cette chaîne ne cherche donc pas à
sélectionner arbitrairement une configuration gagnante, mais à soumettre un
ensemble de variantes à un processus progressif d'évaluation, de qualification,
de sélection et d'agrégation.

La documentation interne du projet formalise cette logique sous la forme d'un
pipeline en sept couches, notées de A à G. Cette structuration répond à trois
problèmes classiques de la recherche quantitative : le biais de sélection,
l'instabilité temporelle et la redondance informationnelle. Elle permet aussi de
maintenir une continuité entre méthode statistique et restitution produit.

\section{Fondement théorique : pourquoi un pipeline et non un meilleur indicateur}
Choisir directement la meilleure variante observée sur un historique est une
procédure fragile. Si l'on teste un grand nombre de variantes, certaines
obtiendront de bons résultats par simple hasard. C'est précisément l'un des
enseignements majeurs de la littérature sur le \emph{data snooping} et le
sur-ajustement \cite{white2000,bailey2014}. Dans le contexte présent, chaque
famille produit un univers de variantes, dont certaines peuvent sembler
performantes sans pour autant révéler une vraie structure informative.

L'idée centrale est donc la suivante : plutôt que de retenir un « meilleur »
indicateur, il faut éliminer successivement les variantes manifestement faibles,
puis agréger les survivantes non redondantes. Cette logique rejoint l'esprit des
méthodes d'ensemble, pour lesquelles la combinaison de plusieurs estimateurs
raisonnablement bons est souvent plus robuste qu'un estimateur isolé
\cite{breiman1996}.

\section{Taxonomie des familles et lecture économique}
Dans le périmètre effectivement exploité par l'interface initiale, quatre
familles sont utilisées : \texttt{sma}, \texttt{macd}, \texttt{rsi} et
\texttt{obv}. Elles sont choisies non pour multiplier artificiellement les
signaux, mais parce qu'elles couvrent des dimensions économiques distinctes
\cite{murphy1999,elder1993}.

\begin{table}[H]
\centering
\caption{Familles retenues et rôle économique associé}
\label{tab:signal-families}
\begin{tabular}{p{2.2cm}p{3cm}p{4.8cm}p{4.2cm}}
\toprule
\textbf{Famille} & \textbf{Catégorie} & \textbf{Question traitée} & \textbf{Exemple de règle} \\
\midrule
\texttt{sma} & Tendance & Le prix évolue-t-il au-dessus ou au-dessous de son niveau moyen ? & \(s_t=\mathrm{sign}(C_t-\mathrm{SMA}_n(t))\) \\
\texttt{macd} & Momentum & L'accélération de la tendance devient-elle favorable ou défavorable ? & \(s_t=\mathrm{sign}(\mathrm{MACD}_t-\mathrm{Signal}_t)\) \\
\texttt{rsi} & Oscillation & Le marché est-il dans une zone d'excès susceptible de se corriger ? & \(s_t=+1\) si \(\mathrm{RSI}_t<\theta_{bas}\) \\
\texttt{obv} & Volume & Les flux de volume confirment-ils le mouvement observé sur les prix ? & \(s_t=\mathrm{sign}(\mathrm{OBV}_t-\mathrm{EMA}(\mathrm{OBV})_t)\) \\
\bottomrule
\end{tabular}
\end{table}

Cette diversité est importante : une famille de tendance et une famille
oscillatoire ne racontent pas la même chose. Le score final ne représente donc
pas une répétition de la même idée sous plusieurs formes, mais une combinaison
raisonnée de perspectives différentes sur un même titre.

\section{Le rôle structurant de l'horizon}
L'horizon n'est pas un simple paramètre d'affichage. Il conditionne à la fois la
géométrie des fenêtres d'évaluation, la profondeur historique exploitée et les
grilles paramétriques testées au sein de chaque famille. Le projet adopte les
conventions suivantes :
\begin{itemize}
\item \textbf{court terme} : thèse de quelques jours à quelques semaines, avec
1 an de contexte et des fenêtres de test trimestrielles ;
\item \textbf{moyen terme} : thèse de plusieurs semaines à plusieurs mois, avec
2 ans de contexte et des fenêtres semestrielles ;
\item \textbf{long terme} : thèse structurelle de plusieurs mois à un an et plus,
avec 3 ans de contexte et des fenêtres annuelles.
\end{itemize}

D'un point de vue mathématique, la configuration par horizon peut être résumée
par le quadruplet
\[
H = (L_{train},L_{test},L_{step},Y_{max})
\]
où \(L_{train}\) est la longueur de mise en contexte, \(L_{test}\) la longueur
de test hors échantillon, \(L_{step}\) l'avance entre deux fenêtres, et
\(Y_{max}\) la profondeur historique maximale exprimée en années.

Le recours à des grilles de paramètres différentes selon l'horizon est lui aussi
fondamental. Pour une moyenne mobile, le paramètre est directement lié à une
mémoire temporelle. Pour un oscillateur borné comme le RSI, en revanche,
augmenter indéfiniment la période finit par comprimer l'indicateur autour de sa
zone centrale ; il faut alors ajuster simultanément la période et les seuils de
surachat/survente. Cette attention portée au calibrage par horizon constitue
l'un des éléments de soin méthodologique les plus importants du projet.

\subsection{Grilles candidates effectivement utilisées pour les quatre familles livrées}
Le soin apporté aux horizons apparaît concrètement dans les grilles de
paramètres. Pour les quatre familles effectivement restituées dans l'interface,
les bandes explorées peuvent être résumées comme suit.

\begin{table}[H]
\centering
\caption{Grilles paramétriques des familles déployées}
\label{tab:implemented-grids}
\begin{tabular}{p{2cm}p{3.2cm}p{3.3cm}p{3.3cm}}
\toprule
\textbf{Famille} & \textbf{Court terme} & \textbf{Moyen terme} & \textbf{Long terme} \\
\midrule
\texttt{sma} & fenêtres rapides, de 3 à 90 barres & fenêtres intermédiaires, de 10 à 250 barres & fenêtres structurelles, de 20 à 400 barres \\
\texttt{macd} & couples rapides autour de \((6\!-\!15,16\!-\!26,7\!-\!9)\) & voisinage du canonique \((12,26,9)\) avec lenteurs jusqu'à 30 & couples lents allant jusqu'à \((20,35,12)\) \\
\texttt{rsi} & périodes 5 à 15, seuils \((30,70)\), \((25,75)\), \((20,80)\) & périodes 7 à 30, mêmes seuils & périodes 10 à 40, mêmes seuils avec lecture plus rare \\
\texttt{obv} & EMA de référence sur la même bande que \texttt{sma} court terme & EMA de référence sur la même bande que \texttt{sma} moyen terme & EMA de référence sur la même bande que \texttt{sma} long terme \\
\bottomrule
\end{tabular}
\end{table}

Le Tableau~\ref{tab:implemented-grids} montre bien que la question n'est pas
seulement de « tester beaucoup de paramètres », mais de tester des paramètres
cohérents avec une intention temporelle donnée. La même famille n'exprime donc
pas la même hypothèse économique lorsqu'elle est évaluée en court terme ou en
long terme.

\section{Vue d'ensemble du pipeline A-G}
Le pipeline peut se résumer ainsi :
\begin{align*}
\text{A} &: \text{génération des candidats},\\
\text{B} &: \text{évaluation hors échantillon},\\
\text{C} &: \text{scoring de robustesse},\\
\text{D} &: \text{filtrage des survivants},\\
\text{E} &: \text{réduction de redondance},\\
\text{F} &: \text{calcul du signal courant},\\
\text{G} &: \text{agrégation finale en score de famille}.
\end{align*}

La sortie de la couche G constitue ensuite l'entrée logique des interfaces de
restitution. Les sous-sections suivantes détaillent chacune de ces couches.

\section{Couche A : génération structurée de l'univers candidat}
Pour chaque couple \((f,h)\), où \(f\) désigne une famille et \(h\) un horizon,
la couche A génère un ensemble déterministe de variantes :
\[
\mathcal{V}_{f,h} = \{v_1, v_2, \dots, v_{30}\}.
\]

Le choix d'un ensemble de 30 variantes par famille répond à un compromis entre
couverture paramétrique, coût computationnel et maîtrise du risque de faux
positifs. Une grille trop étroite conduirait à une exploration insuffisante,
tandis qu'une grille exhaustive augmenterait le coût du \emph{multiple testing}
\cite{harvey2016}.

Pour les familles simples, comme \texttt{sma} ou \texttt{obv}, chaque variante
est essentiellement définie par une fenêtre ou une période de lissage. Pour les
familles multi-paramètres, comme \texttt{macd}, la variante est un triplet de la
forme
\[
v = (p_{fast},p_{slow},p_{sig}), \qquad p_{fast}<p_{slow}.
\]

Les grilles ne sont pas arbitraires. Elles sont choisies de manière à respecter
la signification économique des horizons. Ainsi, un \texttt{MACD}(12,26,9)
appartient naturellement à une lecture de moyen terme, alors qu'une variante
beaucoup plus rapide sera plus adaptée au court terme. Cette calibration est
explicitement documentée dans la couche signal et constitue un garde-fou contre
les choix paramétriques opportunistes.

\section{Couche B : évaluation hors échantillon par fenêtres roulantes}
La couche B est le coeur empirique de la méthode. Pour chaque variante
\(v \in \mathcal{V}_{f,h}\), on calcule un tableau de signaux causaux
\(s_t(v) \in \{-1,0,+1\}\), puis on l'évalue sur une suite de fenêtres
hors échantillon.

Soit une série de clôtures \((C_t)_{t=1}^{T}\). Pour un horizon donné, les
fenêtres sont construites en respectant l'ordre chronologique :
\[
W_k = [t_k,\, t_k+L_{train}) \cup [t_k+L_{train},\, t_k+L_{train}+L_{test}).
\]

La première partie de \(W_k\) fournit le contexte de calcul ; la seconde seule
sert à l'évaluation. Le pas d'avancement vaut \(L_{step}=L_{test}\), ce qui
évite tout recouvrement entre les fenêtres de test. Cette géométrie suit les
recommandations de l'analyse Walk-Forward pour les séries temporelles
\cite{pardo2008,tashman2000}.

Pour une fenêtre valide, les rendements journaliers sont calculés par
\[
r_{t+1}=\frac{C_{t+1}}{C_t}-1.
\]
Si \(p_t\) désigne la position effective induite par le signal au temps \(t\),
le rendement net s'écrit
\begin{equation}
R^{net}_{t+1} = p_t\,r_{t+1} - \kappa\,|p_t-p_{t-1}|,
\label{eq:net-return}
\end{equation}
où \(\kappa=\texttt{cost\_bps}/10^4\) représente le coût proportionnel de
transaction. L'équation~\ref{eq:net-return} traduit un point essentiel : la
méthode ne mesure pas une qualité abstraite de signal, mais une qualité
évaluée \emph{après} coût de rotation. Une variante très instable est donc
naturellement pénalisée.

\subsection{Différence entre signaux d'état et signaux d'action}
Toutes les familles ne produisent pas la même sémantique immédiate. Certaines,
comme \texttt{sma} ou \texttt{obv}, expriment directement un \emph{état} de
position : le marché est lu comme haussier, baissier ou neutre. D'autres, comme
\texttt{rsi}, expriment plutôt une \emph{intention d'action} à partir d'un excès
de marché. Dans ce second cas, le signal élémentaire est converti en position
effective avant évaluation, selon une logique de maintien jusqu'au prochain
événement significatif.

Si \(a_t\in\{-1,0,+1\}\) désigne un signal d'action, la position utilisée dans
l'évaluation peut être représentée de manière récursive par
\begin{equation}
p_t =
\begin{cases}
a_t, & \text{si } a_t \neq 0,\\
p_{t-1}, & \text{sinon.}
\end{cases}
\label{eq:actions-to-positions}
\end{equation}

L'équation~\ref{eq:actions-to-positions} rappelle qu'un oscillateur n'est pas
interprété comme une simple suite de signaux indépendants, mais comme une
instruction qui produit un état de position jusqu'à la prochaine révision.

Sur chaque fenêtre, plusieurs métriques sont produites. Le ratio de Sharpe
annualisé s'écrit
\begin{equation}
\mathrm{Sharpe}_k = \frac{\mu_k}{\sigma_k}\sqrt{252},
\label{eq:sharpe-window}
\end{equation}
où \(\mu_k\) et \(\sigma_k\) désignent respectivement la moyenne et l'écart
type des rendements nets dans la fenêtre \(k\). On calcule également le rendement
cumulé, le nombre de changements de position et le \emph{maximum drawdown}. Une
fenêtre n'est conservée que si elle contient au moins 20 barres effectivement
évaluables.

\section{Couche C : scoring de robustesse multi-critère}
Une variante ne peut pas être classée sur son seul Sharpe moyen. La couche C
synthétise donc l'ensemble des résultats hors échantillon dans un score de
robustesse composite. Avant cela, un filtre de viabilité est appliqué : une
variante n'est étudiée plus loin que si elle possède au moins trois fenêtres
valides et si la fraction de fenêtres à Sharpe positif atteint un seuil minimal.

Formellement, si \(n_{valid}\) désigne le nombre de fenêtres valides et
\(\pi^{+}\) la fraction de fenêtres dont le Sharpe est positif, la viabilité
s'écrit
\begin{equation}
\mathrm{Viable}(v)=\mathbf{1}\{n_{valid}\geq 3\}\,\mathbf{1}\{\pi^{+}\geq 0.40\}.
\label{eq:viability-gate}
\end{equation}

Pour les variantes viables, quatre composantes normalisées sont ensuite calculées.
\begin{align}
S_{sh}(v) &= \mathrm{clip}\!\left(\frac{\overline{Sh}(v)+0.5}{2},0,1\right), \\
S_{st}(v) &= \pi^{+}(v), \\
S_{co}(v) &= \mathrm{clip}\!\left(1-\frac{\sigma_{Sh}(v)}{|\overline{Sh}(v)|+0.1},0,1\right), \\
S_{dd}(v) &= \mathrm{clip}\!\left(1-\frac{\overline{DD}(v)}{0.30},0,1\right).
\end{align}

Le score final de robustesse est alors donné par
\begin{equation}
\mathcal{R}(v)=0.35\,S_{sh}(v)+0.30\,S_{st}(v)+0.20\,S_{co}(v)+0.15\,S_{dd}(v).
\label{eq:robustness-score}
\end{equation}

L'équation~\ref{eq:robustness-score} donne la priorité à la qualité de
performance hors échantillon et à la stabilité temporelle, tout en conservant un
contrôle explicite sur la cohérence et sur la profondeur des pertes. Cette
construction réduit le risque de classement artificiel par un unique indicateur,
ce qui rejoint les critiques formulées contre les approches mono-métriques
\cite{bailey2014}.

\section{Couche D : filtrage des survivants compétitifs}
La couche D vise à réduire l'ensemble des variantes scorées à un sous-ensemble de
survivants véritablement compétitifs. Cette réduction se fait en deux temps.

D'abord, on impose un plancher absolu :
\[
\mathcal{V}^{floor}_{f,h}=\{v\in \mathcal{V}_{f,h} : \mathcal{R}(v)\geq 0.25\}.
\]
Ce plancher empêche de retenir la « moins mauvaise » variante d'une cohorte
faible.

Ensuite, les variantes restantes sont ordonnées par score décroissant et l'on
conserve la fraction supérieure de la cohorte. Si \(q=0.40\) désigne la part
éliminée, l'ensemble des survivants peut se noter
\[
\mathcal{V}^{surv}_{f,h}=\mathrm{Top}_{1-q}(\mathcal{V}^{floor}_{f,h}).
\]

La logique de cette couche n'est pas de trancher définitivement entre les
meilleures variantes, mais de préparer l'étape suivante en éliminant les cas
clairement insuffisants. Le pipeline accepte donc une certaine pluralité de
survivants, à condition qu'ils soient de qualité minimale.

\section{Couche E : réduction de redondance par corrélation}
Même après filtrage, plusieurs survivants peuvent porter une information quasi
identique. C'est le cas, par exemple, de variantes proches d'une même famille de
moyennes mobiles. La couche E élimine ce double comptage informationnel.

Pour chaque survivant \(v\), on considère son vecteur de signaux causaux
\(\mathbf{s}(v)\). On calcule ensuite la corrélation de Pearson entre paires de
survivants :
\begin{equation}
\rho(v_i,v_j)=\frac{\mathrm{Cov}(\mathbf{s}(v_i),\mathbf{s}(v_j))}{\sigma(\mathbf{s}(v_i))\sigma(\mathbf{s}(v_j))}.
\label{eq:pearson-signals}
\end{equation}

L'algorithme est glouton. Les survivants sont triés par score décroissant. On
sélectionne d'abord la variante la plus fiable, puis chaque variante suivante
n'est retenue que si sa corrélation absolue avec tous les représentants déjà
choisis reste sous un seuil fixé. Formellement, une variante \(v\) est admise
dans l'ensemble des représentants \(\mathcal{V}^{rep}_{f,h}\) si
\begin{equation}
\max_{u\in\mathcal{V}^{rep}_{f,h}} |\rho(v,u)| \leq 0.85.
\label{eq:redundancy-threshold}
\end{equation}

Cette étape matérialise un principe simple : l'ensemble final doit combiner des
signaux suffisamment bons \emph{et} suffisamment distincts. Elle rapproche la
construction d'ensemble d'une logique de diversification informationnelle,
proche dans son esprit de la diversification en portefeuille \cite{markowitz1952}.

\section{Couche F : calcul du signal courant des représentants}
Une fois les représentants choisis, la couche F calcule pour chacun d'eux le
signal courant, c'est-à-dire le verdict associé à la dernière barre disponible.
Si \(v\) est un représentant et \(T\) la date la plus récente, on extrait
simplement
\[
s^{cur}(v)=s_T(v) \in \{-1,0,+1\}.
\]

Cette couche est distincte de la couche B pour une raison importante : la couche
B juge la qualité historique du signal, tandis que la couche F interroge son
état actuel. Une variante peut avoir une bonne robustesse historique et émettre,
au dernier instant, un signal neutre ; inversement, une variante active au temps
présent ne sera jamais utilisée si sa robustesse passée est insuffisante.

À ce niveau, chaque représentant est enrichi par une étiquette textuelle adaptée
à sa catégorie : \emph{haussier / baissier}, \emph{survendu / suracheté},
\emph{accumulation / distribution}, etc. Le projet ne réduit donc pas toutes les
familles à une simple sémantique « achat / vente » ; il préserve la lecture
économique propre à chaque type d'indicateur.

\section{Couche G : agrégation finale en score de famille}
La dernière couche combine les représentants en un score unique de famille. Si
\(\mathcal{V}^{rep}_{f,h}=\{v_1,\dots,v_m\}\), le score brut est donné par
\begin{equation}
\mathcal{S}_{f,h}=\frac{\sum_{i=1}^{m} \mathcal{R}(v_i)\,s^{cur}(v_i)}{\sum_{i=1}^{m} \mathcal{R}(v_i)}.
\label{eq:family-ensemble-raw}
\end{equation}

Le score affiché dans l'interface est ensuite exprimé en pourcentage sur
l'intervalle \([-100,100]\) :
\begin{equation}
\mathrm{Score}_{f,h}^{(\%)} = 100\times \mathcal{S}_{f,h}.
\label{eq:family-ensemble-pct}
\end{equation}

Les équations~\ref{eq:family-ensemble-raw} et
\ref{eq:family-ensemble-pct} montrent deux choses. Premièrement, les
représentants les plus robustes ont mécaniquement plus de poids dans la
synthèse. Deuxièmement, un score proche de zéro n'est pas un échec du système :
il traduit simplement une absence de consensus suffisant entre les représentants.

L'étiquette finale de famille résulte alors d'un découpage de l'axe des scores
en zones sémantiques. Pour une famille de tendance, par exemple, une zone
fortement positive sera interprétée comme \emph{très haussière}, une zone proche
de zéro comme \emph{neutre}, et une zone fortement négative comme
\emph{très baissière}. Pour une famille oscillatoire, le vocabulaire associé est
naturellement différent.

\section{Discussion sur la cohérence d'ensemble}
Le pipeline A-G présente plusieurs propriétés importantes du point de vue d'un
mémoire d'ingénierie quantitative.
\begin{itemize}
\item Il est \textbf{causal} : le calcul des signaux et l'évaluation respectent
strictement l'ordre temporel.
\item Il est \textbf{modulaire} : l'ajout d'une nouvelle famille ne modifie pas
la structure des couches C à G.
\item Il est \textbf{horizon-aware} : les paramètres testés et la géométrie des
fenêtres changent avec la thèse temporelle étudiée.
\item Il est \textbf{explicable} : chaque score final peut être redescendu vers
les représentants et vers les filtres qui ont éliminé les autres variantes.
\end{itemize}

Cette dernière propriété est particulièrement importante dans le cadre du projet.
La page Signaux et le tableau de bord ne sont pas de simples écrans de sortie ;
ils prolongent directement la philosophie du pipeline en exposant les résultats
selon le même ordre logique : consensus, familles, représentants.
